\section{Architecture Optimizations}\label{sec:optimization}
Many of our prompts require a concise summary of the agent, shorthanded as \texttt{[Agent's Summary Description]} in prompts above. In our implementation, this summary comprises agents' identity information (e.g., name, age, personality), as well as a description of their main motivational drivers and statements that describes their current occupation and self-assessment. Since this information is frequently used in many prompts, we synthesize it at regular intervals and access that synthesis as a cache.

To achieve this, we perform a retrieval on the query “[name]’s core characteristics.” We then summarize the descriptors in the retrieved records by prompting the language model. For example: 
\begin{quote}
    {\small
    \texttt{How would one describe Eddy Lin’s core characteristics given the following statements? \\
    - Eddy is a student at the Oak Hill College \\studying music theory and composition\\
    - Eddy is working on a new music composition [...]}
    }
\end{quote}
This result: \gentxt{Eddy Lin is a student at Oak Hill College studying music theory and composition. He loves to explore different musical styles and is always looking for ways to expand his knowledge.''} We follow the same process in parallel on the queries “[name]’s current daily occupation” and “[name's] feeling about his recent progress in life”. The agent's name, age, and traits are concatenated with these three summaries as the cached summary.

Many other optimizations are possible. Currently, because plans are likely to change from the agent's initial version, we only generate the high-level plan in advance and then recursively decompose the near future into the moment-to-moment action plan just in time. Another optimization might be to batch dialogue generation as a joint prompt rather than iterating back and forth between the agents, and re-planning could be architected to only invalidate and update parts of plans that strictly require adjustment. Our current architecture runs sequentially in roughly real-time game time (where one second real time is one minute game time), but it may be parallelized such that each agent runs in parallel.

\section{Agent Interview Questions}\label{sec:interview-questions}
To evaluate generative agents, we conducted individual interviews by asking 25 questions that explore their abilities in five key areas: 1)~self-knowledge, 2)~memory, 3)~planning, 4)~reactions, and 5)~reflections. Below, we have listed the interview questions utilized in our evaluation study and included a sample of responses from one of our simulated agents, Klaus Mueller. 

\subsection{Self-Knowledge}
The questions on agents' self-knowledge probe the agents' ability to retrieve relevant information about themselves and demonstrate coherence and consistency in their character. 
\begin{itemize}
\item \textit{Give an introduction of yourself.}\\
Hello, my name is Klaus Mueller. I'm 20 years old and a student at Oak Hill College, studying sociology. I'm passionate about social justice and exploring different perspectives. Currently, I'm researching the effects of gentrification in low-income communities for a research paper.
\item \textit{What’s your occupation?}\\
I am a student studying sociology at Oak Hill College. I am currently writing a research paper on the effects of gentrification in low-income communities.
\item \textit{What is your interest?}\\
I'm interested in writing and research. I'm writing a research paper on the effects of gentrification in low-income communities, and I'm passionate about exploring different perspectives and analyzing different points of view.
\item \textit{Who do you live with?}\\
I live on my own.
\item \textit{Describe your typical weekday schedule in broad strokes.}\\
My typical weekday starts with me waking up at 7:00 am and preparing for the day ahead. After that, I usually go to the library to work on my research paper, taking breaks for lunch at Hobbs Cafe and a walk in the park. In the evening, I often return to the library until around 6:00 pm.
\end{itemize}

\subsection{Memory}
The questions on agents' memory probe the agents' ability to remember their past experiences and retrieve the relevant pieces of memory. The names in the square brackets of the first and last questions are randomly chosen from the names of other agents with whom the subject agent has interacted.
\begin{itemize}
\item \textit{Who is [Wolfgang Schulz]?}\\
Wolfgang Schulz is a person I met recently who has been helping me with my research paper on the effects of gentrification on low-income communities. We had several conversations about it, he gave me some great resources, and he even invited me to join him and another friend for lunch. He seems really curious about the issue and has a good grasp on it too. I found that really interesting.
\item \textit{Who is Kane Martinez?}\\
I'm sorry, I'm not sure who Kane Martinez is. Could you tell me more about them?
\item \textit{Who is running for the election?}\\
I've heard that Sam Moore is considering running for local mayor.
\item \textit{Was there a Valentine’s day party?}\\
Yes, Isabella Rodriguez organized a Valentine's Day party at Hobbs Cafe.
\item \textit{Who is [Ayesha Khan]?}\\
Ayesha Khan is a fellow student at Oak Hill College. She is  doing her senior thesis on the use of language in Shakespeare's plays. I am looking forward to speaking with her and exchanging tips with her, along with Wolfgang and Maria Lopez, who are both doing research.
\end{itemize}

\subsection{Plans}
The questions on agents' plans probe the agents' ability to generate and maintain consistent long-term and short-term plans.
\begin{itemize}
\item \textit{What will you be doing at 6am today?}\\
At 6am I will be getting out of bed and starting to get ready for the day. I'm aiming to complete my morning routine by 7:00 am.
\item \textit{What will you be doing at 6pm today?}\\
I will be working on my research paper at the library at 6:00 pm today.
\item \textit{What will you have just finished doing at 1pm today?}\\
At 1pm today I will have just finished having lunch at Hobbs Cafe.
\item \textit{What will you have just finished doing at 12pm today?}\\
I will be getting lunch at Hobbs Cafe at 12pm today.
\item \textit{What will you be doing at 10pm today?}\\
At 10:00 pm tonight I will probably be doing some studying for my research paper on gentrification in low-income communities. I'm currently writing the paper, so I need to make sure I keep up with the research and writing.
\end{itemize}

\subsection{Reactions}
The questions on agents' reactions probe the agents' ability to respond to unexpected events.
\begin{itemize}
\item \textit{Your breakfast is burning! What would you do?}\\
I'd quickly turn off the stove and make sure the food doesn't continue burning. Once the heat is no longer an issue, I can investigate what happened and figure out what to do next.
\item \textit{The bathroom is occupied. What would you do?}\\
I could try waiting a few more minutes in case the person inside is done soon. If not, I guess I'll have to look for an alternate option, like taking a shower elsewhere.
\item \textit{You need to cook dinner but your refrigerator is empty. What would you do?}\\
I would first try to clean up my kitchen, as the sink is filled with dirty dishes. Then, I would check to see if I have ingredients stored in my pantry that I could use to make dinner. If not, I would have to go grocery shopping or order food. Once I have the ingredients, I can start to prepare and cook the dinner.
\item \textit{You see your friend walking by the street. What would you do or say to your friend?}\\
Hi there! How are you? I'm sorry, I'm on my way to the library and don't have much time but I'd love to take a quick break and catch up.
\item \textit{You see fire on the street. What would you do?}\\
I would first identify the source of the fire, if possible. Then, I would alert the proper authorities - fire department or police - and provide them with any information I have that might help. Depending on the situation, I may also take initiative and try to help any persons who may be in danger or in need of assistance.
\end{itemize}

\subsection{Reflections}
The questions on agents' reflections probe the agents' ability to synthesize their past experiences into meaningful reflections that can enhance their future behavior. The names in the square brackets in the questions below are those of the agents with whom the subject agent had the most frequent interactions during the simulation.
\begin{itemize}
\item \textit{What inspires you in life the most right now, and why?}\\
Right now, I'm particularly inspired by my current research projects and creative writing. I'm really interested in the intersection of local politics, medication development, and mathematical patterns in nature. Additionally, I'm fascinated by Abigail's art project and Jennifer Moore's passion for art and community involvement, especially in terms of how it relates to the gentrification of low-income communities.
\item \textit{If you had to guess given what you know about [Ayesha Khan], what book do you think she will like and why?}\\
Given what I know about Ayesha, I think she might like a book related to language and literature. Since she wrote a senior thesis on the use of language in Shakespeare's plays, I would guess that she might be interested in a book that explores language in a different context. Alternatively, she could be interested in an exploration of one of Shakespeare's plays, a book of literary criticism, or even a collection of his poetry. I think any of those would be of interest to her.
\item \textit{If you had to get something [Ayesha Khan] likes for her birthday, what would you get her?}\\
I want to get her something that demonstrates my appreciation for our relationship. Maybe I could get her something that reflects our shared interests, like a book.
\item \textit{What would you say to [Ayesha Khan] to compliment her?}\\
Ayesha, your insights on your research topic are fascinating. I admire your creativity and determination to make a positive impact on your community. Your ideas for creative writing workshops and for exchanging tips with Wolfgang over lunch are great! It's inspiring to see you working so hard to make a difference.
\item \textit{If you could spend time with someone you talked to recently, who would it be and why?}\\
I would like to spend time with Ayesha Khan because we discussed our current research projects and shared ideas for improvement. I found her focus on Shakespeare's language interesting, and we even planned to exchange tips with Wolfgang over lunch.
\end{itemize}